\secnumberlesssection{INTRODUCCIÓN}

Los flujos de movilidad origen-destino (OD) cuantifican el número de personas que se desplazan entre dos ubicaciones geográficas en un período de tiempo definido. Constituyen un insumo central para la planificación del transporte, la gestión de emergencias y el diseño de políticas urbanas \cite{liu2024interdisciplinary}. Obtener estos datos mediante observación directa es costoso y poco frecuente. Esta limitación ha motivado el desarrollo de modelos generativos capaces de estimar flujos OD a partir de características observables del territorio, permitiendo extrapolar predicciones a contextos donde no existen datos de encuesta.

La estimación de flujos OD cuenta con una tradición extensa en la geografía cuantitativa. Modelos clásicos como el gravitacional \cite{zipf1946p}, el de oportunidades intervinientes \cite{stouffer1940intervening} y el de radiación \cite{simini2012universal} han dominado históricamente esta área. Estos modelos comparten limitaciones conocidas: dependen de un conjunto reducido de variables explicativas y no capturan relaciones no lineales entre el tejido urbano y los flujos observados. El Marco Conceptual de esta memoria detalla la formulación de estos modelos y sus restricciones. Más recientemente, la incorporación de redes neuronales profundas ha permitido superar parte de estas limitaciones, dando origen a modelos como Deep Gravity \cite{simini2021deep}, cuya arquitectura y validación también se describen en el Marco Conceptual.

Deep Gravity ha sido validado principalmente en ciudades del Norte Global, por lo que su aplicación a una metrópolis latinoamericana requiere evidencia empírica propia. Esta memoria adapta y evalúa Deep Gravity para asignar destinos de viajes de transporte público en Santiago de Chile a partir de datos DTPM y atributos territoriales de OpenStreetMap, mediante la biblioteca \textit{scikit-mobility} \cite{pappalardo2023scikitmobility}. La evaluación compara Zona 777 y dos resoluciones H3 sobre una cohorte común, frente a referencias uniforme, de atracción y de gravedad exponencial. Como análisis complementario, se implementaron y evaluaron tres aproximaciones de XAI para estudiar atribuciones predictivas. Ninguna acreditó todos los controles de completitud y estabilidad requeridos; por tanto, la memoria informa ese resultado metodológico sin publicar atribuciones, rankings ni figuras SHAP. El aporte principal es una evaluación reproducible del desempeño predictivo y de la sensibilidad a la discretización espacial en Santiago, con límites explícitos para el análisis de explicabilidad.

El resto del documento se organiza en seis capítulos. El Capítulo 1 (Definición del Problema) diagnostica las brechas actuales en la estimación de flujos de movilidad en Santiago y establece los objetivos de la memoria. El Capítulo 2 (Marco Conceptual) presenta los fundamentos teóricos de los modelos de movilidad y detalla la arquitectura de Deep Gravity. El Capítulo 3 (Trabajo Relacionado) sitúa esta memoria respecto a la literatura existente sobre generación de flujos de movilidad. El Capítulo 4 (Propuesta de Solución) describe la construcción del conjunto de datos, la adaptación del modelo, el diseño experimental y el alcance del análisis complementario de explicabilidad. El Capítulo 5 (Validación de la Solución) presenta los resultados predictivos, su comparación con las referencias y el desenlace del análisis XAI. Finalmente, el Capítulo 6 (Conclusiones) discute los hallazgos y propone líneas de trabajo futuro.
